% BANKS-SOURCE: pdf=88 print=78 kind=implicit-exercise id=I-CH07-002
\begin{exercise}[$O(n)$ multiplets and crossing]{I-CH07-002}
Let $\phi^a$, $a=1,\ldots,n$, be real scalar fields in the defining
representation of an unbroken $O(n)$ symmetry.  Assume that their exact
two-point function has an isolated one-particle pole and that the associated
particles are stable.  For the process
$\phi^a(p_1)\phi^b(p_2)\to\phi^c(p_3)\phi^d(p_4)$, take
\[
\begin{aligned}
 p_i^2&=-m^2,\qquad s=-(p_1+p_2)^2,\qquad t=-(p_1-p_3)^2,\\
 u&=-(p_1-p_4)^2,\qquad s+t+u=4m^2.
\end{aligned}
\]
\begin{enumerate}[label=(\alph*)]
\item Use $O(n)$ invariance and the Lehmann representation to determine the
masses and pole residues of the one-particle states.

\item Apply LSZ to determine the most general invariant rank-four tensor
structure of the amplitude.  Use Bose symmetry and crossing to express its
coefficient functions through one scalar function at crossed arguments, and
state the relations under exchange of either particle pair.

\item Decompose the amplitude into the singlet, antisymmetric, and
symmetric-traceless $s$-channel representations.

\item Explain which additional rank-four tensors are allowed when $O(n)$ is
replaced by $SO(n)$, including the low-dimensional cases.
\end{enumerate}
\end{exercise}
